% GENERATED FILE. Edit canonical lessons, then run npm run generate:books.
% canonical-source-hash: fnv1a64:1414af23f35615d1
% canonical-lessons: SA-C427-vasatih, SA-C427-patamandapah, SA-C427-nivesah, SA-C427-astrnati, SA-C427-prabudhyate

\chapter{A Night's Lodging}
\label{ch:sa-a-night-s-lodging}
\hlchaptermodality{427}

\begin{chapteropening}
\textbf{By the end of this chapter:} \emph{I can say where I spent the night and when I woke.}

Complete the last lesson of chapter 427: I can say where I spent the night and when I woke.
\end{chapteropening}

\section[\texorpdfstring{\sk{वसतिः} (vasatiḥ)}{वसतिः (vasatiḥ)}]{\sk{वसतिः} (vasatiḥ) — a night's stay; a lodging for the night}

\label{lesson:SA-C427-vasatih}



\begin{quote}
Before the new one: say the Sanskrit for stays on, then the Sanskrit for a place of refuge.
\end{quote}

\subsection*{You'll want to know: \sk{वसतिः}}
\textbf{\sk{वसतिः}} — \emph{vasatiḥ} — \textquotedblleft{}a night's stay; a lodging for the night\textquotedblright{}.

\begin{quote}
\sk{पथिकः} \sk{ग्रामे} \sk{वसतिं} \sk{याचते।}
\end{quote}

\begin{quote}
\emph{pathikaḥ grāme vasatiṃ yācate.}
\end{quote}

\begin{quote}
The wayfarer asks for a night's stay in the village.
\end{quote}

A feminine noun.

\begin{cousinweb}
\textbf{\sk{वसतिः}} (\emph{vasatiḥ}) is built on the root √vas, \textquotedblleft{}to dwell, to stay\textquotedblright{}, which also gives \emph{vasati}, \textquotedblleft{}dwells\textquotedblright{}; MW gives the noun \emph{vasati} as staying, especially overnight, a dwelling (MW).
\end{cousinweb}

\begin{grammarlens}[title={What you've built}]
One more word for getting around.
\end{grammarlens}

\subsection*{Your turn}
\begin{itemize}
\raggedright
  \item \emph{Say it:} \emph{vasatiḥ}
  \item \emph{Say it:} \emph{vasatiḥ}, once more
  \item \emph{Say it:} \emph{pathikaḥ grāme vasatiṃ yācate.} — The wayfarer asks for a night's stay in the village.
  \item \emph{Recall:} say the Sanskrit for stays on, then the Sanskrit for a place of refuge, then say \emph{vasatiḥ} again
\end{itemize}

\begin{tcolorbox}[breakable,colback=teal!4,colframe=teal!35!black,title={Before you move on}]
What does \sk{वसतिः} mean? (\textquotedblleft{}A night's stay; a lodging for the night\textquotedblright{}.) Say it once more.
\end{tcolorbox}
\section[\texorpdfstring{\sk{पटमण्डपः} (paṭamaṇḍapaḥ)}{पटमण्डपः (paṭamaṇḍapaḥ)}]{\sk{पटमण्डपः} (paṭamaṇḍapaḥ) — a tent}

\label{lesson:SA-C427-patamandapah}



\begin{quote}
Before the new one: say the Sanskrit for a place of refuge, then the Sanskrit for a night's stay.
\end{quote}

\subsection*{You'll want to know: \sk{पटमण्डपः}}
\textbf{\sk{पटमण्डपः}} — \emph{paṭamaṇḍapaḥ} — \textquotedblleft{}a tent\textquotedblright{}.

\begin{quote}
\sk{पर्वते} \sk{वयं} \sk{पटमण्डपे} \sk{स्वपिमः।}
\end{quote}

\begin{quote}
\emph{parvate vayaṃ paṭamaṇḍape svapimaḥ.}
\end{quote}

\begin{quote}
On the mountain we sleep in a tent.
\end{quote}

A masculine noun.

\begin{cousinweb}
\textbf{\sk{पटमण्डपः}} (\emph{paṭamaṇḍapaḥ}) is built from \emph{paṭa}, \textquotedblleft{}cloth, canvas\textquotedblright{}, and \emph{maṇḍapa}, \textquotedblleft{}a pavilion, an open hall\textquotedblright{}; MW gives the compound as a canvas house, a tent (MW).
\end{cousinweb}

\begin{grammarlens}[title={What you've built}]
One more word for getting around.
\end{grammarlens}

\subsection*{Your turn}
\begin{itemize}
\raggedright
  \item \emph{Say it:} \emph{paṭamaṇḍapaḥ}
  \item \emph{Say it:} \emph{paṭamaṇḍapaḥ}, once more
  \item \emph{Say it:} \emph{parvate vayaṃ paṭamaṇḍape svapimaḥ.} — On the mountain we sleep in a tent.
  \item \emph{Recall:} say the Sanskrit for a place of refuge, then the Sanskrit for a night's stay, then say \emph{paṭamaṇḍapaḥ} again
\end{itemize}

\begin{tcolorbox}[breakable,colback=teal!4,colframe=teal!35!black,title={Before you move on}]
What does \sk{पटमण्डपः} mean? (\textquotedblleft{}A tent\textquotedblright{}.) Say it once more.
\end{tcolorbox}
\section[\texorpdfstring{\sk{निवेशः} (niveśaḥ)}{निवेशः (niveśaḥ)}]{\sk{निवेशः} (niveśaḥ) — an encampment; a halting place}

\label{lesson:SA-C427-nivesah}



\begin{quote}
Before the new one: say the Sanskrit for a night's stay, then the Sanskrit for a tent.
\end{quote}

\subsection*{You'll want to know: \sk{निवेशः}}
\textbf{\sk{निवेशः}} — \emph{niveśaḥ} — \textquotedblleft{}an encampment; a halting place\textquotedblright{}.

\begin{quote}
\sk{सार्थस्य} \sk{निवेशः} \sk{नद्याः} \sk{तीरे} \sk{अस्ति।}
\end{quote}

\begin{quote}
\emph{sārthasya niveśaḥ nadyāḥ tīre asti.}
\end{quote}

\begin{quote}
The caravan's encampment is on the riverbank.
\end{quote}

A masculine noun.

\begin{cousinweb}
\textbf{\sk{निवेशः}} (\emph{niveśaḥ}) is built from \emph{ni}, \textquotedblleft{}down, in\textquotedblright{}, and the root √viś, \textquotedblleft{}to enter, to settle\textquotedblright{}; MW gives \emph{niveśa} as settling in a place, encamping, halting, and a dwelling or camp (MW).
\end{cousinweb}

\begin{grammarlens}[title={What you've built}]
One more word for getting around.
\end{grammarlens}

\subsection*{Your turn}
\begin{itemize}
\raggedright
  \item \emph{Say it:} \emph{niveśaḥ}
  \item \emph{Say it:} \emph{niveśaḥ}, once more
  \item \emph{Say it:} \emph{sārthasya niveśaḥ nadyāḥ tīre asti.} — The caravan's encampment is on the riverbank.
  \item \emph{Recall:} say the Sanskrit for a night's stay, then the Sanskrit for a tent, then say \emph{niveśaḥ} again
\end{itemize}

\begin{tcolorbox}[breakable,colback=teal!4,colframe=teal!35!black,title={Before you move on}]
What does \sk{निवेशः} mean? (\textquotedblleft{}An encampment; a halting place\textquotedblright{}.) Say it once more.
\end{tcolorbox}
\section[\texorpdfstring{\sk{आस्तृणाति} (āstṛṇāti)}{आस्तृणाति (āstṛṇāti)}]{\sk{आस्तृणाति} (āstṛṇāti) — spreads out}

\label{lesson:SA-C427-astrnati}



\begin{quote}
Before the new one: say the Sanskrit for a tent, then the Sanskrit for an encampment.
\end{quote}

\subsection*{You'll want to know: \sk{आस्तृणाति}}
\textbf{\sk{आस्तृणाति}} — \emph{āstṛṇāti} — \textquotedblleft{}spreads out\textquotedblright{}.

\begin{quote}
\sk{पथिकः} \sk{भूमौ} \sk{कम्बलम्} \sk{आस्तृणाति।}
\end{quote}

\begin{quote}
\emph{pathikaḥ bhūmau kambalam āstṛṇāti.}
\end{quote}

\begin{quote}
The wayfarer spreads a blanket on the ground.
\end{quote}

Past (he …): \sk{सः} \sk{कम्बलम्} \sk{आस्तीर्णवान्}, \emph{saḥ kambalam āstīrṇavān}.

\begin{cousinweb}
\textbf{\sk{आस्तृणाति}} (\emph{āstṛṇāti}) is built from \emph{ā} and the root √stṛ, \textquotedblleft{}to spread, to strew\textquotedblright{}, which also gives \emph{āstaraṇa}, \textquotedblleft{}a carpet, a spread\textquotedblright{}; MW gives \emph{ā-stṛ} as to spread or strew over, with the participle \emph{āstīrṇa}, \textquotedblleft{}spread\textquotedblright{} (MW).
\end{cousinweb}

\begin{grammarlens}[title={What you've built}]
One more word for getting around.
\end{grammarlens}

\subsection*{Your turn}
\begin{itemize}
\raggedright
  \item \emph{Say it:} \emph{āstṛṇāti}
  \item \emph{Say it:} \emph{āstṛṇāti}, once more
  \item \emph{Say it:} \emph{pathikaḥ bhūmau kambalam āstṛṇāti.} — The wayfarer spreads a blanket on the ground.
  \item \emph{Recall:} say the Sanskrit for a tent, then the Sanskrit for an encampment, then say \emph{āstṛṇāti} again
\end{itemize}

\begin{tcolorbox}[breakable,colback=teal!4,colframe=teal!35!black,title={Before you move on}]
What does \sk{आस्तृणाति} mean? (\textquotedblleft{}Spreads out\textquotedblright{}.) Say it once more.
\end{tcolorbox}
\section[\texorpdfstring{\sk{प्रबुध्यते} (prabudhyate)}{प्रबुध्यते (prabudhyate)}]{\sk{प्रबुध्यते} (prabudhyate) — wakes up}

\label{lesson:SA-C427-prabudhyate}



\begin{quote}
Before the new one: say the Sanskrit for an encampment, then the Sanskrit for spreads out.
\end{quote}

\subsection*{You'll want to know: \sk{प्रबुध्यते}}
\textbf{\sk{प्रबुध्यते}} — \emph{prabudhyate} — \textquotedblleft{}wakes up\textquotedblright{}.

\begin{quote}
\sk{पथिकः} \sk{प्रातः} \sk{प्रबुध्यते।}
\end{quote}

\begin{quote}
\emph{pathikaḥ prātaḥ prabudhyate.}
\end{quote}

\begin{quote}
The wayfarer wakes up in the morning.
\end{quote}

Past (he …): \sk{सः} \sk{प्रातः} \sk{प्रबुद्धवान्}, \emph{saḥ prātaḥ prabuddhavān}.

\begin{cousinweb}
\textbf{\sk{प्रबुध्यते}} (\emph{prabudhyate}) is built from \emph{pra}, \textquotedblleft{}forth\textquotedblright{}, and the root √budh, \textquotedblleft{}to wake, to know\textquotedblright{}, which also gives \emph{budhyate}, \textquotedblleft{}understands\textquotedblright{}; MW gives \emph{pra-budh} as to wake up, with the participle \emph{prabuddha}, \textquotedblleft{}awakened\textquotedblright{} (MW).
\end{cousinweb}

\begin{grammarlens}[title={What you've built}]
One more word for getting around.
\end{grammarlens}

\subsection*{Your turn}
\begin{itemize}
\raggedright
  \item \emph{Say it:} \emph{prabudhyate}
  \item \emph{Say it:} \emph{prabudhyate}, once more
  \item \emph{Say it:} \emph{pathikaḥ prātaḥ prabudhyate.} — The wayfarer wakes up in the morning.
  \item \emph{Recall:} say the Sanskrit for an encampment, then the Sanskrit for spreads out, then say \emph{prabudhyate} again
\end{itemize}

\begin{tcolorbox}[breakable,colback=teal!4,colframe=teal!35!black,title={Before you move on}]
What does \sk{प्रबुध्यते} mean? (\textquotedblleft{}Wakes up\textquotedblright{}.) Say it once more.
\end{tcolorbox}
